\begin{claim}{C.28}{Proved}{Orbit types of a loop-generated plane group}
Let $t\mapsto f_t$ be a compact-open continuous based loop in
$\operatorname{Homeo}^+(\R^2)$, with $f_0=f_1=\id$, and put
\[
 L=\{f_t:0\leq t\leq1\},\qquad P=L\cup L^{-1},
 \qquad G=\langle P\rangle.
\]
Every $G$-orbit is one of the following: an open subset of the plane;
a point; a Jordan curve; or the image of a continuous injective map
$\eta:\R\to\R^2$ that restricts to a homeomorphism on every compact
interval. In the last case its intrinsic arc order identifies it with
the line, but no global homeomorphism from $\R$ to the orbit with its
subspace topology is asserted. The loop restricts continuously in the
compact-open topology to the intrinsically ordered line. Every
homeomorphism commuting with $G$ carries these orbit types to the same
types and induces a homeomorphism between the corresponding intrinsic
lines, allowing reversal of their orders.
\end{claim}
\begin{proof}
Fix $x\in\R^2$ and put
\[
                  K_n=P^n x,\qquad O=Gx=\bigcup_{n\geq1}K_n.
                                                        \tag{C.28.1}
\]
The sets $K_n$ are nested compact Peano continua. Indeed, $P$ is a
continuous image of the wedge of two circles, using the given loop
and its inverse; it is compact, connected and locally connected.
A finite product of such spaces has the same properties, and the
word-evaluation map has image $K_n$. To recall why local connectedness
passes through this compact surjection, for an open neighborhood of a
point in the image cover its compact inverse fibre by finitely many
connected open subsets of its inverse image. Their images all meet
that point and lie in its component. The complement of the image of
the remaining compact set is a smaller open neighborhood in that
component. Thus components of open sets in the image are open.
In particular $K_n$ is arcwise connected. The identity lies in $P$,
which proves the nesting in (C.28.1). Every element of $G$ can be
joined to the identity by a path in $G$, by writing it as a finite word
and using the specified paths to its factors.

Suppose first that $O$ contains a Jordan curve $J$. For any $g\in G$
choose a path $g_s$ from $\id$ to $g$, so that $g_sJ\subseteq O$.
Let $D_s$ be the bounded Jordan domain of $g_sJ$. If $gJ\ne J$, then
$D_0\mathbin\triangle D_1$ contains a nonempty open set away from the
two boundary curves. Every point $q$ in this open set lies on some
$g_sJ$: otherwise the winding number of the continuously varying curve
$g_sJ$ around $q$ would be constant, whereas its endpoint values are
$0$ and $1$. Hence $O$ contains a nonempty open set. Transitivity of
$G$ on $O$ carries this open set to a neighborhood of each point of
$O$, so $O$ is open. If every $g$ preserves $J$, then for a point
$y\in J$ its orbit is contained in $J$; since $Gy=O$ and $J\subseteq O$,
we obtain $O=J$. Thus a nonopen orbit containing a Jordan curve is
itself that Jordan curve.

We may now consider a nonopen orbit containing no Jordan curve.
Each $K_n$ is a dendrite: a compact, connected, locally connected
metric space containing no Jordan curve. Arcs between two points are
unique, since two distinct arcs would contain a Jordan curve. This
also proves uniqueness of arcs in all of $O$: two proposed arcs in
$O$ would give the same forbidden curve. Existence follows by putting
their endpoints into one $K_n$.

Each $K_n$ has countably many branch points. Choose a countable dense
subset $D_n$. Three distinct components of the complement of a branch
point are open there, so choose a point of $D_n$ in each. The branch
point is the unique common point of the three pairwise connecting
arcs for this triple: two branches meet only at their point of
separation. Thus each branch point is determined by one of the
countably many triples from $D_n$.

A triod point of $O$ is a branch point of some $K_n$. Put the three
outer endpoints of its triod into one $K_n$; the unique connecting
arcs in $K_n$ must be the arcs of that triod. The triod points of $O$
are consequently countable. Their set is invariant under every element
of $G$. If it were nonempty, transitivity would make every point of
$O$ a triod point. A nontrivial arcwise connected orbit contains an
arc and is uncountable, giving a contradiction. Thus a nontrivial
$O$ has no triod points, and each nondegenerate $K_n$ is an arc.

The nested arcs give $O$ a consistent linear arc order, unique up to
reversal. To see the absence of additional branches directly, fix two
distinct points $a,b$. Every other point is on $[a,b]$, beyond $a$,
or beyond $b$, because any fourth alternative would give a triod.
There are at most two endpoints of this ordered union. The endpoint
set is invariant under $G$, so transitivity and uncountability exclude
endpoints. We can therefore pass to a cofinal sequence of arcs $A_j$
such that $A_j$ lies in the relative interior of $A_{j+1}$: extend
both endpoints in $O$ and then put the two extensions and the next
$K_n$ into a later $K_m$.

Choose compatible homeomorphisms $[-j,j]\to A_j$, extending at each
step along the two added arc intervals. Their union is a continuous
injective map
\[
                 \eta:\R\longrightarrow\R^2,
                        \qquad\eta(\R)=O,              \tag{C.28.2}
\]
whose restriction to every compact interval is a homeomorphism onto
its image. These assertions follow on each finite interval from one
of the compact arc parametrizations. This supplies the stated
intrinsic line structure. It does not prove that the inverse of
(C.28.2) is continuous for the global subspace topology; accumulations
from far along an orbit have not been excluded.

The restricted loop is continuous for that intrinsic line topology.
A compact arc in $O$ is contained in some $K_m$, and
$f_t(K_m)\subseteq K_{m+1}$ for all $t$. The latter is a compact arc,
on which the intrinsic and subspace topologies coincide. Joint
continuity of $(t,z)\mapsto f_t(z)$ on each such compact arc therefore
gives compact-open continuity of
$\eta^{-1}f_t\eta\in\operatorname{Homeo}(\R)$.
Its order sign is positive throughout the loop: images of two distinct
points vary continuously in the line and cannot cross. Finally a
commuting homeomorphism $h$ satisfies the exact equality
\[
                  h(K_n(x))=K_n(hx).
\]
It therefore maps compact arcs homeomorphically to compact arcs and
induces a homeomorphism of the intrinsic ordered lines; its inverse
has the same property. It also preserves the point, Jordan-curve and
open alternatives, completing the proof.
\end{proof}

\paragraph{Exact application and remaining step.}
The canonical lift in the paragraph after C.12 sends a hyperbolic
centralizer loop to a closed based loop in
$\operatorname{Homeo}^+(\widetilde U)$, where
$\widetilde U\cong\R^2$. The entire group $\widetilde H$ of lifts of
the operator image commutes with it and hence with the $G$ of C.28.
Thus C.28 applies to this actual lift group, and C.27 provides an
order-natural contraction on each intrinsic line orbit after
restriction to its stabilizer in $\widetilde H$.
These are typed restriction and lift maps, but they do not yet give a
homotopy through homeomorphisms of the whole plane. Continuity across
distinct orbit strata, compatibility with their frontiers, and a
relative contraction on proper open orbits remain unproved.
C.26 treats global transitivity on a closed hyperbolic surface; a
proper open orbit in C.28 does not satisfy that hypothesis. The
admissible infinite-group level-$0$ question therefore remains C.12.
